\documentclass[12pt]{article}
\usepackage[margin=1in]{geometry}
\usepackage{amsmath}
\usepackage{mathtools}
\usepackage{amssymb}
\usepackage{lmodern}
\usepackage{cancel}
\usepackage{tikz}
\usepackage{graphicx}
\setlength{\parskip}{0.3em}
\setlength{\parindent}{0pt}
% 使用系统字体，需要 XeLaTeX 或 LuaLaTeX 编译
\usepackage{fontspec}
\setmainfont{Times New Roman}
\usepackage{xcolor}
\usepackage{float}
\usepackage{caption}
\usepackage{indentfirst}
\setlength{\parindent}{2em}
\newcommand{\mdavg}[2]{\langle #1 \rangle\!\langle #2 \rangle}
\newcommand{\avg}[1]{\langle #1 \rangle}
\newcommand{\doubleavg}[2]{\left\langle #1 \right\rangle\!\left\langle #2 \right\rangle}
\newcommand{\inavg}[2]{\langle #1 \rangle\! [#2]}
\newcommand{\rinavg}[2]{[#1]\!\langle #2 \rangle}
\newcommand{\aket}[1]{|#1\rangle}
\newcommand{\asqu}[1]{{\langle#1\rangle}^2}
\newcommand{\sket}[1]{|#1]}
\newcommand{\cbrak}[2]{\avg{#1}\![#2]}
\newcommand{\acbrak}[2]{[#1]\!\avg{#2}}
\newcommand{\abra}[1]{\langle #1|}
\newcommand{\sbra}[1]{[#1|}
\newcommand{\inner}[2]{\langle #1 | #2 \rangle}
\newcommand{\smallkern}{\kern0.08em}
\begin{document}

\title{Some Details}
\author{Su Yingze}
\date{\today}
\maketitle
\section{Normalization of Polarization Vectors}
First, we give the definition of polarization vectors which is written in terms of adjoint representation
\begin{equation}
    \varepsilon_a^\mu=\frac{\abra{p_I}(r\cdot \Gamma)\Gamma^\mu\aket{p^J}(\sigma_a)^I_{\ J}}{4(p\cdot r)}
\end{equation}
here $a$ runs from 1 to 3. We can use antisymmetric tensor to rewrite this equation,like 
\begin{equation}
    \varepsilon_a^\mu=\frac{\abra{p_I}(r\cdot \Gamma)\Gamma^\mu\aket{p_K}\epsilon^{JK}(\sigma_a)^I_{\ J}}{4(p\cdot r)}
\end{equation}
here we choose the left-multiplication convention.

Then, a $2\times 2$ symmetric matrix can be defined like 
\begin{equation}
    (\sigma_a)^{IJ}=\epsilon^{JK}(\sigma_a)^I_{\ K}.
\end{equation}
This means that the original polarization vector can be understood as 
\begin{align}
    \varepsilon_a^\mu&=-\frac{\abra{p_I}(r\cdot \Gamma)\Gamma^\mu\aket{p_J}}{4(p\cdot r)}(\sigma_a)^{IJ}\notag \\
    &=-\varepsilon^\mu_{IJ}(\sigma_a)^{IJ}
\end{align}
here $I,J$ are fundamental little group indices, runs from 1 to 2.

And we know that the reasonable polarization vector for physical state should be symmetric with little group indices, it means that we need 
to symmetrize the indices $I$ and $J$ as $(IJ)=\frac{1}{2}(IJ+JI)$. Because of the symmetry of $(\sigma_a)^{IJ}$, so it can be obtained that
\begin{equation}
    \varepsilon_a^\mu =-(\varepsilon^\mu_{(IJ)}+\varepsilon^\mu_{[IJ]})(\sigma_a)^{IJ}=-\varepsilon^\mu_{(IJ)}(\sigma_a)^{IJ}
\end{equation}
which means we only need to consider the symmetric part of $\varepsilon^\mu_{IJ}$.

Now, we want to compute the normalization for $\varepsilon^\mu_{IJ}$, and after multiplying the symmetric matrix, it should go back to the normalization
\begin{equation}
    \varepsilon_a^\mu\varepsilon_{b,\mu}=-\delta_{ab}
\end{equation}

The method to compute the normalization is just multiplying two polarization vectors directly, like
\begin{equation}
    \varepsilon^\mu_{(IJ)}\varepsilon^{}_{(KL),\mu}=\frac{\abra{p_I}(r\cdot \Gamma)\Gamma^\mu\aket{p_J}}{4(p\cdot r)}\times \frac{\abra{p_K}(r\cdot \Gamma)\Gamma_\mu\aket{p_L}}{4(p\cdot r)}
    \label{7}
\end{equation}
here we need to use a special 5d equality
\begin{equation}
    (\Gamma^\mu)_A^{\ B}(\Gamma_\mu)_C^{\ D}=-2\Omega_{AC}\Omega^{BD}+2\delta_A^{\ D}\delta_C^{\ B}-\delta_A^{\ B}\delta_C^{\ D}
    \label{8}
\end{equation}
and the dirac operator can be decomposed to
\begin{equation}
    \cancel{p}=p\cdot \Gamma=p_A^{\ B}= \aket{p_I}_A\abra{p^I}^B
\end{equation}
which help us to rewrite the equation \eqref{7} to be
\begin{align}
    \varepsilon^\mu_{IJ}\varepsilon^{}_{KL,\mu}&=\frac{\abra{p_I}^A\aket{r_M}_A\abra{r^M}^B(\Gamma^\mu)_B^{\ C}\aket{p_J}_C \abra{p_K}^D\aket{r_N}_D\abra{r^N}^E(\Gamma_\mu)_D^{\ F}\aket{p_L}_F}{16(p\cdot r)^2}\notag \\
    &=\frac{\inner{p_I}{r_M}\inner{p_K}{r_N}}{16(p\cdot r)^2}\times \abra{r^M}^B(\Gamma^\mu)_B^{\ C}\aket{p_J}_C \abra{r^N}^E(\Gamma_\mu)_E^{\ F}\aket{p_L}_F
    \label{10}
\end{align}
After plugging the equation \eqref{8} into the second part of \eqref{10}, it becomes
\begin{align}
    &\quad \abra{r^M}^B \aket{p_J}_C \abra{r^N}^E \aket{p_L}_F (-2\Omega_{BE}\Omega^{CF}+2\delta_B^{\ F}\delta_E^{\ C}-\delta_B^{\ C}\delta_E^{\ F})\notag \\
    &=-2\inner{r^M}{r^N}\inner{p_L}{p_J}+2\inner{r^M}{p_L}\inner{r^N}{p_J}-\inner{r^M}{p_J}\inner{r^N}{p_L}
    \label{11}
\end{align}

The first term of \eqref{11} equals to 0 because the equality
\begin{equation}
    \inner{p_I}{p_J}=0
\end{equation}
and the remaining two term can be contracted with the first part of \eqref{10}, which helps to obtain the following two terms
\begin{equation}
    2\inner{p_I}{r_M}\inner{r^M}{p_L}\inner{p_K}{r_N}\inner{r^N}{p_J}-\inner{p_I}{r_M}\inner{r^M}{p_J}\inner{p_K}{r_N}\inner{r^N}{p_L}
\end{equation}
It is known that 
\begin{equation}
    \abra{p_I}(r\cdot \Gamma) \aket{p_J}=\mathrm{Tr}[\aket{p_J}\abra{p_I}(r\cdot \Gamma)]=\frac{1}{2}\epsilon_{IJ}\mathrm{Tr}[(p\cdot \Gamma)(r\cdot \Gamma)]=2\epsilon_{IJ}(p\cdot r)
\end{equation}
so the formal equation can be rewritten like 
\begin{equation}
    8(p\cdot r)^2\epsilon_{IL}\epsilon_{KJ}-4(p\cdot r)^2\epsilon_{IJ}\epsilon_{KL}
\end{equation}

Finally, we need to symmetrize the indices $(IJ)$ and $(KL)$, so that we obtain
\begin{equation}
    \varepsilon^\mu_{IJ}\varepsilon^{}_{KL,\mu}=\frac{1}{8(p\cdot r)^2}(p\cdot r)^2(\epsilon_{IL}\epsilon_{KJ}+\epsilon_{JL}\epsilon_{KI}+\epsilon_{IK}\epsilon_{LJ}+\epsilon_{JK}\epsilon_{LI})-\frac{1}{16 (p\cdot r)^2}(p\cdot r)^2(\epsilon_{(IJ)}\epsilon_{(KL)})
\end{equation}
because $\epsilon$ itself is an antisymmetric tensor so it will become zero under the symmetrization, so that the second term is vanishing.

After some rearrangement of indices, we can obtain
\begin{equation}
    \varepsilon^\mu_{IJ}\varepsilon^{}_{KL,\mu}=\frac{1}{4}(\epsilon_{IL}\epsilon_{KJ}+\epsilon_{IK}\epsilon_{LJ})
\end{equation}

Unfortunately, when we want to reproduce the original normalization relation, we will obtain something like
\begin{align}
    \varepsilon_a^\mu\varepsilon_{b,\mu}&=(-\varepsilon^\mu_{(IJ)}(\sigma_a)^{IJ})(-\varepsilon^\mu_{(KL)}(\sigma_b)^{KL})\notag \\
    &= \frac{1}{4}(\epsilon_{IL}\epsilon_{KJ}+\epsilon_{IK}\epsilon_{LJ})(\sigma_a)^{IJ}(\sigma_b)^{KL}\notag \\
    &=\frac{1}{4}\epsilon_{IL}\epsilon_{KJ}\epsilon^{JM}(\sigma_a)^I_{\ M}\epsilon^{LN}(\sigma_b)^K_{\ N}+\frac{1}{4}\epsilon_{IK}\epsilon_{LJ}\epsilon^{IM}(\sigma_a)^J_{\ M}\epsilon^{LN}(\sigma_b)^K_{\ N}\notag \\
    &=\frac{1}{4}\delta_I^{\ N}\delta_K^{\ M}(\sigma_a)^I_{\ M}(\sigma_b)^K_{\ N}+\frac{1}{4}\delta_K^{\ M}\delta_J^{\ N}(\sigma_a)^J_{\ M}(\sigma_b)^K_{\ N}\notag \\
    &=\frac{1}{4}2\mathrm{tr}(\sigma_a\sigma_b)\notag \\
    &=\delta_{ab}
\end{align}

\section{Completeness Relation of Polarization Vectors}
For this purpose, we need to compute the following formula
\begin{equation}
    \varepsilon_a^\mu \varepsilon^{a,\nu}=?
\end{equation}
If we multiply the momentum $p_\mu$, it should become vanishing so that the completeness relation should look like 
\begin{equation}
    \varepsilon_a^\mu \varepsilon^{a,\nu}=k(-\eta^{\mu\nu}+\frac{p^\mu r^\nu+p^\nu r^\mu}{(p\cdot r)^2})
\end{equation}
here k is just a overall coefficient.

Here, we choose to compute the completeness relation for polarization vectors with fundamental indices
\begin{equation}
    \varepsilon_{IJ}^\mu \varepsilon^{IJ,\nu}=a(-\eta^{\mu\nu}+\frac{p^\mu r^\nu+p^\nu r^\mu}{(p\cdot r)^2})
\end{equation}
with a different coefficient.

For symplisity, we start with the convention from paper arxiv:2202.08257, in which the polarization vector looks like
\begin{equation}
    \epsilon_{IJ}^\mu=\frac{\abra{p_{(I}}\Gamma^\mu \aket{r_{J)}}}{\sqrt{2}}
\end{equation}
with the normalization
\begin{equation}
    \inner{p^I}{r^J}=\epsilon^{IJ}
    \label{22}
\end{equation}
which means the following condition
\begin{equation}
    2(p\cdot r)=1
\end{equation}
The method to compute this completeness relation is still just multiplying them togther like
\begin{equation}
    \varepsilon_{(IJ)}^\mu \varepsilon^{(IJ),\nu}=\frac{\abra{p_{(I}}\Gamma^\mu\aket{r_{J)}}}{\sqrt{2}}\times \frac{\abra{p^{(I}}\Gamma^\nu\aket{r^{J)}}}{\sqrt{2}}
\end{equation}

It can be decomposed into 4 pices $IJ^{\textstyle IJ}$, $IJ^{\textstyle JI}$, $JI^{\textstyle IJ}$, $JI^{\textstyle JI}$, and the first one can be computed like
\begin{align}
    \abra{p_I}\Gamma^\mu\aket{r_J}\times \abra {p^I}\Gamma^\nu\aket{r^{J}}&=-\abra{p_I}\Gamma^\mu\aket{r_J}\abra {r^J}\Gamma^\nu\aket{p^{I}}\notag\\
    &=-\mathrm{Tr}[\aket{p^{I}}\abra{p_I}\Gamma^\mu\aket{r_J}\abra r^J\Gamma^\nu]\notag \\
    &=\mathrm{Tr}[(p\cdot \Gamma)\Gamma^\mu (r\cdot \Gamma)\Gamma^\nu]\notag \\
    &=4(p^\mu r^\nu+p^\nu r^\mu -(p\cdot r)\eta^{\mu\nu})
\end{align}
where we use a nontrival fact that 
\begin{equation}
    \abra {p^I}\Gamma^\nu\aket{r^{J}}=-\abra {r^J}\Gamma^\nu\aket{p^{I}}\quad (\text{obtained from}\quad C^T\Gamma^\mu C=-(\Gamma^\mu)^T)
\end{equation}
The fourth piece is nearly the same as the first one
\begin{equation}
    \abra{p_J}\Gamma^\mu\aket{r_I}\times \abra {p^J}\Gamma^\nu\aket{r^{I}}=4(p^\mu r^\nu+p^\nu r^\mu -(p\cdot r)\eta^{\mu\nu})
\end{equation}
The remaining two need some extra efforts, here I compute it explictly
\begin{equation}
    \abra{p_I}\Gamma^\mu\aket{r_J}\times \abra {p^J}\Gamma^\nu\aket{r^{I}}
    \label{27} 
\end{equation}
In terms of two massless spinors that satisfy \eqref{22}, the $USp(2,2)$ identity operator can be written as
\begin{equation}
    \aket{r_a}_A\abra{p^a}^B+\aket{p_a}_A\abra{r^a}^B=\delta_A^{\ B}
\end{equation}
It can be used to rewrite the equation \eqref{27} like
\begin{align}
    &\quad \ \abra{p_I}^A(\Gamma^\mu)_A^{\ B}\aket{r_J}_B\abra{p^J}^C(\Gamma^\nu)_C^{\ D}\aket{r^I}_D \notag \\
    &=\abra{p_I}^A(\Gamma^\mu)_A^{\ B} (\delta_B^{\ C}-\aket{p_J}_B\abra{r^J}^C)(\Gamma^\nu)_C^{\ D}\aket{r^I}_D \notag \\
    &=\abra{p_I}^A(\Gamma^\mu)_A^{\ B}(\Gamma^\nu)_B^{\ D}\aket{r^I}_D-\abra{p_I}\Gamma^\mu\aket{p_J}\abra{r^J}\Gamma^\nu\aket{r^I}
\end{align}
And similarly for 
\begin{align}
    \abra{p_J}\Gamma^\mu\aket{r_I}\times \abra {p^I}\Gamma^\nu\aket{r^{J}}&=\abra {p^I}\Gamma^\nu\aket{r^{J}}\times \abra{p_J}\Gamma^\mu\aket{r_I}\notag \\
    &=\abra {p^I}^A(\Gamma^\nu)_A^{\ B}(-\aket{p^J}_B\abra{r_J}^{\ C}-\delta_B^{\ C})(\Gamma^\mu)_C^{\ D}\aket{r_I}_D \notag \\
    &=-\abra {p^I}^A(\Gamma^\nu)_A^{\ B}(\Gamma^\mu)_B^{\ D}\aket{r_I}_D-\abra{p^I}\Gamma^\nu\aket{p^J}\abra{r_J}\Gamma^\mu\aket{r_I}
\end{align}
Now we can combine these two terms
\begin{align}
    &\abra{p_I}^A(\Gamma^\mu)_A^{\ B}(\Gamma^\nu)_B^{\ D}\aket{r^I}_D-\abra{p_I}\Gamma^\mu\aket{p_J}\abra{r^J}\Gamma^\nu\aket{r^I}-\abra {p^I}^A(\Gamma^\nu)_A^{\ B}(\Gamma^\mu)_B^{\ D}\aket{r_I}_D-\abra{p^I}\Gamma^\mu\aket{p^J}\abra{r_J}\Gamma^\nu\aket{r_I}\notag \\
    &=\abra{p_I}(\Gamma^\mu\Gamma^\nu+\Gamma^\nu\Gamma^\mu)\aket{r^I}-\abra{p_I}\Gamma^\mu\aket{p_J}\abra{r^J}\Gamma^\nu\aket{r^I}-\abra{p^I}\Gamma^\nu\aket{p^J}\abra{r_J}\Gamma^\mu\aket{r_I}\notag \\
    &=2\eta^{\mu\nu}\inner{p_I}{r^I}-(2p^\mu\epsilon_{IJ})(-2r^\nu\epsilon^{JI})-(-2p^\nu\epsilon^{IJ})(2r^\mu\epsilon_{JI})\notag \\
    &=4\eta^{\mu\nu}+8p^\mu r^\nu+8p^\nu r^\mu
\end{align}

Above all, we can conclude that 
\begin{align}
    \varepsilon_{(IJ)}^\mu \varepsilon^{(IJ),\nu}&=\frac{1}{2}\times \frac{1}{8}[8(p^\mu r^\nu+p^\nu r^\mu -(p\cdot r)\eta^{\mu\nu})+4\eta^{\mu\nu}+8p^\mu r^\nu+8p^\nu r^\mu]\notag \\
    &\overset{?}{=}2p^{(\mu}r^{\nu)}
\end{align}

\section{Completeness Relation Review}
Under the following definition of polarization vector
\begin{equation}
    \varepsilon_a^\mu=\frac{\abra{p_I}(r\cdot \Gamma)\Gamma^\mu\aket{p_J}}{4(p\cdot r)}(\sigma_a)^I_{\ K}\epsilon^{KJ}
\end{equation}
and same for 
\begin{equation}
    \varepsilon_a^\mu=\frac{\abra{p^I}(r\cdot \Gamma)\Gamma^\mu\aket{p^J}}{4(p\cdot r)}(\sigma_a)^K_{\ J}\epsilon_{KI},
\end{equation}
the polarization vector living in the adjoint representation can be defined as
\begin{equation}
    \varepsilon^\mu_{IJ}(p;r)\coloneqq \frac{\abra{p_{(I}}(r\cdot \Gamma)\Gamma^\mu\aket{p_{J)}}}{4(p\cdot r)}
\end{equation}
\begin{equation}
    \varepsilon^{IJ,\mu}(p;r)\coloneqq \frac{\abra{p^{(I}}(r\cdot \Gamma)\Gamma^\mu\aket{p^{J)}}}{4(p\cdot r)}.
\end{equation}

The corresponding complex conjugate can be obtained
\begin{equation}
    (\varepsilon^\mu_{IJ}(p;r))^*=-\frac{\abra{p^{(I}}\Gamma^\mu(r\cdot \Gamma)\aket{p^{J)}}}{4(p\cdot r)}
\end{equation}
and similarly
\begin{equation}
    (\varepsilon^{IJ,\mu}(p;r))^*=-\frac{\abra{p_{(I}}\Gamma^\mu(r\cdot \Gamma)\aket{p_{J)}}}{4(p\cdot r)}.
\end{equation}

For the consistency of index placement with the summation convention, we can define the following quantities
\begin{equation}
    \varepsilon^{*\mu}_{IJ}(p;r)\coloneqq  (\varepsilon^{IJ,\mu}(p;r))^*=-\frac{\abra{p_{(I}}\Gamma^\mu(r\cdot \Gamma)\aket{p_{J)}}}{4(p\cdot r)},
\end{equation}
\begin{equation}
    \varepsilon^{*IJ,\mu}(p;r)\coloneqq (\varepsilon^\mu_{IJ}(p;r))^*=-\frac{\abra{p^{(I}}\Gamma^\mu(r\cdot \Gamma)\aket{p^{J)}}}{4(p\cdot r)}
\end{equation}

Then the completeness relation can be computed like
\begin{equation}
    \varepsilon^{*\,IJ,\mu}\,\varepsilon_{IJ}^{\nu}=-\frac{\abra{p^{(I}}\Gamma^\mu(r\cdot \Gamma)\aket{p^{J)}}}{4(p\cdot r)}\frac{\abra{p_{(I}}(r\cdot \Gamma)\Gamma^\nu\aket{p_{J)}}}{4(p\cdot r)}
\end{equation}
It includes 4 terms and can be grouped into two part, the first part can be computed like
\begin{align*}
    &\quad \abra{p^I}\Gamma^\mu(r\cdot \Gamma)\aket{p^J}\abra{p_I}(r\cdot \Gamma)\Gamma^\nu \aket{p_J}\\
    &=-\abra{p^I}\Gamma^\mu(r\cdot \Gamma)\aket{p^J}\abra{p_J}\Gamma^\nu(r\cdot \Gamma) \aket{p_I}\\
    &=-\mathrm{tr}[\aket{p_I}\abra{p^I}\Gamma^\mu(r\cdot \Gamma)\aket{p^J}\abra{p_J}\Gamma^\nu(r\cdot\Gamma)]\\
    &=\mathrm{tr}[(p\cdot \Gamma)\Gamma^\mu(r\cdot \Gamma)(p\cdot \Gamma)\Gamma^\nu (r\cdot \Gamma)] \label{44}\tag{\theequation}\stepcounter{equation}
\end{align*}
Here we need to compute trace for the product of 6 gamma matrices, and I just give the answer
\begin{align*}
\mathrm{Tr}\!\big(\Gamma^{a}\Gamma^{b}\Gamma^{c}\Gamma^{d}\Gamma^{e}\Gamma^{f}\big)=4\Big(
& \eta^{ab}\eta^{cd}\eta^{ef}
-\eta^{ab}\eta^{ce}\eta^{df}
+\eta^{ab}\eta^{cf}\eta^{de} \\
&-\eta^{ac}\eta^{bd}\eta^{ef}
+\eta^{ac}\eta^{be}\eta^{df}
-\eta^{ac}\eta^{bf}\eta^{de} \\
&+\eta^{ad}\eta^{bc}\eta^{ef}
-\eta^{ad}\eta^{be}\eta^{cf}
+\eta^{ad}\eta^{bf}\eta^{ce} \\
&-\eta^{ae}\eta^{bc}\eta^{df}
+\eta^{ae}\eta^{bd}\eta^{cf}
-\eta^{ae}\eta^{bf}\eta^{cd} \\
&+\eta^{af}\eta^{bc}\eta^{de}
-\eta^{af}\eta^{bd}\eta^{ce}
+\eta^{af}\eta^{be}\eta^{cd}
\Big). \tag{\theequation}\stepcounter{equation}
\end{align*}
A compact way to write the result is
\[
\mathrm{Tr}\!\big(\Gamma^{a_1}\cdots\Gamma^{a_6}\big)
= 4 \sum_{\text{pairings}} (-1)^{\sigma}\,
\eta^{a_{i_1} a_{j_1}}\eta^{a_{i_2} a_{j_2}}\eta^{a_{i_3} a_{j_3}} .
\]
Here the sum runs over all pairwise partitions (pairings) of the index set
$\{1,\dots,6\}$ into three disjoint pairs 
$\{(i_1,j_1),(i_2,j_2),(i_3,j_3)\}$.Here $(-1)^{\sigma}$ denotes the sign (parity) of the permutation needed to reorder the sequence so that the members of each pair become adjacent.

Using this formula, we can rewrite the equation \eqref{44}
\begin{align}
    \eqref{44}=&\quad p^\mu r^\nu (p\cdot r)-p^\mu r^\nu (p\cdot r)+p^\mu p^\nu r^2 \notag \\
    &-p^\mu r^\nu(p\cdot r)+\eta^{\mu\nu}(p\cdot r)^2-p^\nu r^\mu (p\cdot r)\notag \\
    &+r^\mu r^\nu p^2-\eta^{\mu\nu} p^2 r^2+r^\mu r^\nu p^2\notag \\
    &-p^\nu r^\mu (p\cdot r)+p^\nu p^\mu r^2 -p^\nu r^\mu (p\cdot r)\notag \\
    &+p^\nu r^\mu (p\cdot r)-p^\mu r^\nu (p\cdot r)+ \eta^{\mu\nu}(p\cdot r)^2\notag \\
    &=4(-2p^\mu r^\nu (p\cdot r)-2p^\nu r^\mu (p\cdot r)+\eta^{\mu\nu}(p\cdot r)^2)
\end{align}
in which many terms cancel with each other.

Similarly, the second part can be computed like
\begin{align}
    &\quad \abra{p^I}\Gamma^\mu(r\cdot \Gamma)\aket{p^J}\abra{p_J}(r\cdot \Gamma)\Gamma^\nu \aket{p_I}\notag \\
    &=\mathrm{tr}[\aket{p_I}\abra{p^I}\Gamma^\mu(r\cdot \Gamma)\aket{p^J}\abra{p_J}(r\cdot\Gamma)\Gamma^\nu]\notag\\
    &=-\mathrm{tr}[(p\cdot \Gamma)\Gamma^\mu(r\cdot \Gamma)(p\cdot \Gamma) (r\cdot \Gamma)\Gamma^\nu]\notag \\
    &=-\mathrm{tr}[(p\cdot \Gamma)\Gamma^\mu(r\cdot \Gamma)(p\cdot \Gamma) \{(r\cdot \Gamma),\Gamma^\nu\}]+\mathrm{tr}[(p\cdot \Gamma)\Gamma^\mu(r\cdot \Gamma)(p\cdot \Gamma)\Gamma^\nu (r\cdot \Gamma)]
\end{align}
It can be noticed that the first term  equals to 0 and the second term is just the same as the result for first part. In short, the two part give us the same result.

Above all. the completeness relation can be written as
\begin{align}
    \varepsilon^{*\,IJ,\mu}\,\varepsilon_{IJ}^{\nu}&=-\frac{1}{16(p\cdot r)^2}\times\frac{1}{4}\times \left[4\times4(-2p^\mu r^\nu (p\cdot r)-2p^\nu r^\mu (p\cdot r)+\eta^{\mu\nu}(p\cdot r)^2)\right]\notag \\
    &=-\frac{1}{2}\left[\frac{-p^\mu r^\nu-p^\nu r^\mu}{(p\cdot r)}+\eta^{\mu\nu}\right]\notag \\
    &=-\frac{1}{2}\eta^{\mu\nu}+\frac{p^{(\mu}r^{\nu)}}{(p\cdot r)}
\end{align}
We can quickly check its correctness by multiplying a momentum
\begin{equation}
    p_\mu\varepsilon^{*\,IJ,\mu}\,\varepsilon_{IJ}^{\nu}=-\frac{1}{2}p^\nu+\frac{\frac{1}{2}p^2 r^\nu+\frac{1}{2}p^\nu(p\cdot r) }{(p \cdot r)}=0
\end{equation}

\section{All for 5 dimension}
\subsection{Spinor-Helicty in 5D}
The Lorentz group in 5D is $\textrm{Spin}(5)\cong \mathrm{USp}(4)\cong \mathrm{Sp}(4,\mathbb{C} )\cap \mathrm{SU}(4) $,we can also consider the non-compact version $\mathrm{SO}(1,4)$, which has an isomorphism $\mathrm{SO}(1,4)\cong \mathrm{USp}(2,2)$. For a 5-momentum vector $p^\mu$, we can identify it with a $4\times 4$ antisymmetric matrix
\begin{equation}
    p_{AB}=p_\mu\gamma^\mu_{AB}\equiv -p_\mu (\gamma^\mu)_A^{\ C}\Omega_{CB}=-(\cancel{p}\ \Omega)_{AB}
    \label{4.1}
\end{equation}
Here we choose a symplectic form to be the metric
\begin{equation}
    \Omega_{AB}=\begin{pmatrix}
        \epsilon_{\alpha\beta} & 0 \\
        0 & -\epsilon^{\dot{\alpha}\dot{\beta}}
    \end{pmatrix}= \Omega^{BA}
\end{equation}
where we use the fact that fundamental indices $A,B$ for $\mathrm{USp}(2,2)$ can be decomposed to $\mathrm{SL}(2,\mathbb{C})$ indices like $A=\alpha\oplus \dot{\alpha} $, and the normalization for 2 dimension Levi-Civita tensor is $\epsilon^{12}=\epsilon_{21}=1$.

We use metric $\Omega$ to lower and raise the indices like 
\begin{equation}
    X^A=\Omega^{AB}X_B,\qquad X_A=\Omega_{AB}X^B
\end{equation} 
where $X$ are spinors belongs to $\mathrm{USp}(2,2)$.

Also, from the defination \eqref{4.1}, we know the gamma matrices with lower indices can be defined like 
\begin{equation}
    \gamma_{AB}=\Omega_{BC}\gamma_A^{\ C}
\end{equation}
similary the gamma matrix with upper indices can be obtained by
\begin{equation}
    \gamma^{AB}=\Omega^{AC}\gamma_C^{\ B}.
\end{equation}
The gamma matrices with lower indices are antisymmetric and $\Omega$-traceless
\begin{equation}
    \gamma_{AB}=\gamma_{BA},\qquad \gamma_{AB}\Omega^{AB}=0
\end{equation}

The determinant of $p_{AB}$ evaluates to $\mathrm{Det}(p_{AB})=\frac{1}{4}(p_{AB}p^{AB})^2=(p^2)^2$. For 5D massless particle, $p^2=0 \Rightarrow \mathrm{Det}(p_{AB}=0)$, so the rank of this matrix equals to 2. Beacuse the little group of 5D massless particle is $\mathrm{SO}(3)\cong \mathrm{SU}(2)$, so the matrix $p_{AB}$ can be decomposed to be the outer product of two $4\times 2$ matrix, here specifically, the $\mathrm{USp}(2,2)$ spinors with $\mathrm{SU}(2)$ little group indices
\begin{equation}
    p_{AB}=\lambda_A^{\ a}\lambda_{Ba}=\aket{p^a}_A\aket{p_a}_B,\qquad p_A^{\ B}=\lambda_A^{\ a}\lambda^B_{\ a}=\aket{p^a}_A\abra{p_a}^B
\end{equation}
with $a,b$ runs over $1,2$. A little group transformation that preserves $p_{AB}$ behaves like $\aket{p^a}_A \rightarrow t^a_{\ b} \aket{p^b}_A$, here $t^a_{\ b}$ refers to the $\mathrm{SU}(2)$ little group transformation. 

Since the little groups of massive 4D momentum and massless 5D momentum are identified, there exist a consistent embedding of the former into the latter.
\begin{equation}
    \lambda_A^{\ a} = 
    \begin{pmatrix}
        \lambda_\alpha^{\;a} \\
        \tilde{\lambda}^{\dot{\alpha} a}
    \end{pmatrix},
    \qquad
    \Omega_{AB} = 
    \begin{pmatrix}
        -\epsilon_{\alpha\beta} & 0 \\
        0 & \epsilon^{\dot{\alpha}\dot{\beta}}
    \end{pmatrix},
\end{equation}
where we can think of 5D Lorentz indices breaking up into 4D Lorentz indices 
\( A \mapsto (\alpha, \dot{\alpha}) \) once a particular direction is chosen, means we manifestly break the 5D Lorentz symmetry. 

The antisymmetry of \( p_{AB} \) implies the important identity that 
\begin{equation}
    \Omega^{AB} \lambda_A^{\;a} \lambda_B^{\;b} = 0,
\end{equation}
so no Lorentz-invariant little-group tensors can be formed from just 
considering a single particle alone.

\subsection{5D 3-particle Amplitude}
\noindent
\textbf{Special Kinematics in 5D}

Scattering amplitudes, are Lorentz invariant and transform covariantly under the little group. So like the massive 4D case, the amplitudes here are $\mathrm{SU}(2)$ covariants, means it carries $\mathrm{SU}(2)$ little group indices and no free Lorentz indices. We want keep the consistency with 4D so here the spin-s particle still transformed in the (2s+1)-dimension symmetric representation of 2s product of $\mathrm{SU}(2)$. In short, a spin-s particle should carry $2s$ $\mathrm{SU}(2)$ little group indices. 

Then, a 3-particle amplitude of spin $s_1,s_2,s_3$ is a tensor carrying $(2s_1+2s_2+2s_3)$ indices like
\begin{equation}
    \mathcal{M}_{(a_1 \cdots a_{2s_1})(b_1 \cdots b_{2s_2})(c_1 \cdots c_{2s_3})}
\end{equation} 
for $a_i, b_j, c_k$ indices of $\mathrm{SU}(2)$ little group of particle 1, 2, 3 respectively.

In order to construct amplitudes which is a tensor carrying little group indices, it is necessary to figure out what tensors structure are possible here. The spinors that make up an on-shell 3-point amplitude without any Lorentz indices can be possibly listed as follows
\begin{gather}
    y^{12}{}_{ab}\coloneqq (\lambda_1)^A\kern0.08em_a (\lambda_2)_{Ab}\\
    y^{23}{}_{bc}\coloneqq (\lambda_2)^A\smallkern_b (\lambda_3)_{Ac}\\
    y^{31}{}_{ca}\coloneqq (\lambda_3)^A\smallkern_c (\lambda_3)_{Aa}
\end{gather}
Because of momentum conservation, the determinant of this $y^{ij}$ vanishes 
\begin{equation}
    \mathrm{Det}(y^{12})=\frac{1}{2}\epsilon^{ab}\epsilon^{cd}y^{12}{}_{ac}y^{12}{}_{bd}=\frac{1}{2}y^{12}{}_{ac}y^{12ac}=\frac{1}{2}p_1^{AB}p_{2AB}=2(p_1\cdot p_2)=0
\end{equation}
This means we can decompose each of them into a outer product of $\mathrm{SU}(2)$ spinors:
\begin{gather}
    y^{12}{}_{ab}=y^1{}_ay^2{}_b\\
    y^{23}{}_{bc}=y^2{}_by^3{}_c\\
    y^{31}{}_{ca}=y^3{}_cy^1{}_a
\end{gather}
In the following, we use the notation: $1_a\coloneqq y^1{}_a$ and $1^A{}_a\coloneqq (\lambda_1)^A{}_a$ for massless 5D spniors.

First, we can try to write down a 3-point amplitude using these variables, e.g. the Yukawa interaction in 5D is precisely $y^12{}_{ab}$:
\begin{equation}
    \mathcal{M}_{a,b}=\text{Corresponding Feynman diagram} = g1_a2_b 
\end{equation}
\textbf{Connecting to 4D $\mathrm{x}$-factor}

If there exist a similar "x" factor in 5D, it would be of the form $\chi^c{}_{c'}$ with no Lorentz indices but little group indices. It can be understood as
\begin{equation}
    p_2^{AB}3_B{}^c=\chi^c{}_{c'}3^{Ac'}
\end{equation}
There is no inverse of $\chi$ because of $\chi^2=0$, $\mathrm{tr} \chi =0$ and $\mathrm{Det}\chi = 0$. In the following, it can be known that 
\begin{equation}
    \chi^c{}_{c'}=3^c3_{c'}
\end{equation}
so the special kinematics $\mathrm{SU}(2)$ spinors $y^3{}_c=3_c$ is a sort of  ``square root" of $\chi$, the x-factor in 5D. Similarly, for other two particles, we can obtain the corresponding x-factor
\begin{gather}
    p_3^{AB}1_B{}^a=\chi^a{}_{a'}1^{Aa'}=(1_a1^{a'})1^{Aa'}\\
    p_1^{AB}2_B{}^b=\chi^b{}_{b'}2^{Ab'}=(2_b2^{b'})2^{Ab'}
\end{gather}

Then, in scalr-QED in 5D, the scalar-scalar-photon amplitude can be directly written as
\begin{equation}
    \mathcal{M}_{cc'}=(\text{Corresponding Feynman diagram})=g3_c3_{c'}
\end{equation}

If one compactifies to 4D, all 4D massive amplitudes should satisfy $m_1\pm m_2 \pm^{'}m_3=0$. If one particle 3 is massless in 4D and other 2 are massive with same mass m, the 4D $x$ factor appears in the form as
\begin{equation}
    3_c3_{c'}=\begin{pmatrix}
        -mx^{-1} & \mp im \\
        \mp im & mx
    \end{pmatrix}_{cc'}
\end{equation}
So we obtain the corresponding amplitude after compactifying, where the diagonal components refer to the massless 4-vector of particle 3, and the off-diagonal components represent the scalar sector. Either $\mathcal{M} \sim  x$ or $x^{-1}$ depending on its helicity.

\appendix
\section{Detailed explaination of special 5D spinor variables}
\textbf{Lemma} : \textit{Momentum consevation for 3-particle amplitude implies}
\begin{gather}
    y^{12}{}_{ab}\coloneqq 1^A{}_a2_{Ab}=y^1{}_ay^2_b=1_a2_b\\
    y^{23}{}_{bc}\coloneqq 2^A{}_b3_{Ac}=y^2{}_by^3_c=2_b3_c\\
    y^{31}{}_{ca}\coloneqq 3^A{}_c1_{Aa}=y^3{}_cy^1_a=3_c1_a
\end{gather}

\noindent
\textit{Proof}. First note that $y^{(i)(i)}_{ab}=0$, e.g.
\begin{equation}
    \Omega_{AB}1^A_a2^B_b=C\epsilon_{ab} (\text{from antisymmetry of ab})\quad \text{and} \quad C \propto \Omega_{AB}1^A_a1^B_b\epsilon^{ab}\propto \Omega_{AB}p_1^{AB}=0
\end{equation}

Another property of y is that the determinant of this $y^{ij}$ vanishes such that 
\begin{equation}
    \mathrm{Det}(y^{12})=\frac{1}{2}\epsilon^{ab}\epsilon^{cd}y^{12}{}_{ac}y^{12}{}_{bd}=\frac{1}{2}y^{12}{}_{ac}y^{12ac}=\frac{1}{2}p_1^{AB}p_{2AB}=2(p_1\cdot p_2)=0
\end{equation}
Thus it is a rank 1 quantity, which can be decomposed into the outer product of 2-spinors:
\begin{align}
y^{12}{}_{ab} &= y^{1}{}_{a} y^{2}{}_{b} \\
y^{23}{}_{bc} &= \Upsilon^{2}{}_{b} \Upsilon^{3}{}_{c} \\
y^{31}{}_{ca} &= \Upsilon^{3}{}_{c} \Upsilon^{1}{}_{a}
\end{align}
Note that not all of there $\mathrm{SU}(2)$ spinors are independent, indeed it is only a 2-dimension vector space. By momentum conservation
\begin{equation}
    y^{12}{}_a{}^b y^{23}{}_b{}^c=1^A{}_a2_A{}^b2^B{}_b3_B{}^c=1^A{}_ap_{aA}{}^B3_B{}^c=1^A{}_a(-p_{1A}{}^B-p_{3A}{}^B)3_B{}^c=0
\end{equation}
Thus $\Upsilon^{2b}y^2{}_b=0$,so $\Upsilon^2{}_b\propto y^2{}_b$ (note the metric we are using now is antisymmetric tensor $\epsilon_{ab}$). Similarly, 
\begin{gather}
    y^{23}{}_b{}^c y^{31}{}_c{}^a=0 \Rightarrow \Upsilon^3{}_c\propto y^3{}_c\\
    y^{31}{}_c{}^a y^{12}{}_a{}^b=0 \Rightarrow \Upsilon^1{}_a\propto y^1{}_a
\end{gather}
Therefore, up to some constants A,B,C, we have 
\begin{align}
    y^{12}{}_{ab}&=y^1{}_ay^2_b\\
    y^{23}{}_{bc}&=By^2{}_by^3_c\\
    y^{31}{}_{ca}&=Cy^3{}_cAy^1_a
\end{align}
As long as we are working over $\mathbb{C} $, we can always rescale them to obtain
\begin{align}
    y^{12}{}_{ab}&=y^1{}_ay^2_b\\
    y^{23}{}_{bc}&=y^2{}_by^3_c\\
    y^{31}{}_{ca}&=y^3{}_cy^1_a
\end{align}
here each y are determined uniquely, up to an overall flip $(y^1, y^2, y^3)\mapsto (-y^1, -y^2, -y^3)$.

The space of 2-dimension representation of $\mathrm{SU}(2)$ is too small to ensure the hold of following equality
\begin{equation}
    \delta_A \coloneqq 1^a1_{Aa}=2^b2_{Ab}=3^c3_{Ac}
\end{equation}
$\delta_A$ is same for all particles.

\textit{Proof}. From the 3-particle momentum conservation
\begin{align}
    0&=1_A{}^a1_{Ba}+2_A{}^b2_{Bb}+3_A{}^c3_{Bc}\\
    0&=\underbrace{1^A{}_a(1_A{}^{a'}}_{=0}1_{B{a'}}+2_A{}^b2_{Bb}+3_A{}^c3_{Bc})\\
    0&=(1_a2^b)2_{Bb}+(-1_a3^c)3_{Bc}\\
    0&=1_a(2^b2_{Bb}-3^c3_{Bc})\\
    &\Rightarrow 2^b2_{Bb}=3^c3_{Bc}
\end{align}
Cyclically premutating the indices, we can obtain $1^a1_{Ba}=2^b2_{Bb}=3^c3_{Bc}$.

A useful equality can be obtained 
\begin{equation}
    \delta_A\lambda_(i)^A{}_b\propto y^{ia}y^i{}_b=0
\end{equation}
\section{$\chi$ in 5D}
Since $\delta_A=1^a1_{Aa}=2^b2_{Ab}=3^c3_{Ac}$, we can evaluate $\chi_3$ by
\begin{equation}
    p_2{}^{AB}3_B{}^c=2^{Ab}2^B{}_b3_B{}^c=2^{Ab}2_b3^c=-\delta^A3^c=3^{A{c'}}3_{c'}3^c=\underbrace{(3^c3_{c'})}_{\chi_3} 3^{A{c'}}
\end{equation}
Cycling the particle labels, we can obtain the $\chi$ for each particle
\begin{gather}
    p_1{}^{AB}2_B{}^b=(2^b2_{b'})2^{A{b'}}\\
    p_3{}^{AB}1_B{}^a=(1^a1_{a'})1^{A{a'}}
\end{gather}
Other combinations can be simply obtained by momentum conservation,like $ p_1{}^{AB}3_B{}^c=-p_2{}^{AB}3_B{}^c=-\chi_3^c{}_{c'}3^{A{c'}}$.
\end{document}